%! Unidad 1 — Leyes fundamentales de la química: ejercicios resueltos y ejercicios finales
%! Fuente: Fuentes/Quimica/Unidad 1 («Ejercicios resueltos», «Practica y calcula», «Piensa y resuelve»).
%! Enunciados adaptados con datos cambiados; soluciones recalculadas. Dificultad: \(\star\) baja, \(\star\star\) media, \(\star\star\star\) alta.

\section{Ejercicios resueltos}\label{sec:ud1-resueltos}

Salvo que se indique otra cosa, usa las masas atómicas H = 1; C = 12; N = 14; O = 16; Na = 23; S = 32; Cl = 35,5;
Fe = 55,8; y \(R = \qty{0.082}{\atm\litre\per\mole\per\kelvin}\).

\begin{example}
    \textbf{Óxidos de hierro.} Se analizan cuatro muestras de óxidos de hierro con estos resultados:
    \begin{itemize}
        \item Muestra A: \qty{8}{\gram} de óxido con un 70\,\% de hierro.
        \item Muestra B: \qty{20}{\gram} de óxido, de los que \qty{14}{\gram} son de hierro.
        \item Muestra C: \qty{4.2}{\gram} de hierro y \qty{1.2}{\gram} de oxígeno.
        \item Muestra D: \qty{2}{\gram} de óxido, de los que \qty{0.6}{\gram} son de oxígeno.
    \end{itemize}
    a) Organiza los datos en una tabla con las masas de hierro, oxígeno y óxido. b) ¿Se trata siempre del mismo óxido?
    c) ¿Qué masa de oxígeno se necesita para que \qty{21}{\gram} de hierro formen un óxido como el de la muestra C?

    \textbf{Solución.} a) Usamos la ley de conservación de la masa, \(m_{\ch{Fe}} + m_{\ch{O}} = m_{\text{óxido}}\):
    A: \(m_{\ch{Fe}} = \num{0.70} \cdot 8 = \qty{5.6}{\gram}\) y \(m_{\ch{O}} = \qty{2.4}{\gram}\); B: \(m_{\ch{O}} = 20 - 14 = \qty{6}{\gram}\);
    C: \(m_{\text{óxido}} = \num{4.2} + \num{1.2} = \qty{5.4}{\gram}\); D: \(m_{\ch{Fe}} = 2 - \num{0.6} = \qty{1.4}{\gram}\).
    \begin{center}
        \begin{tabular}{@{}cccc@{}}
            \toprule
            \textbf{Muestra} & \textbf{Hierro (g)} & \textbf{Oxígeno (g)} & \textbf{Óxido (g)} \\
            \midrule
            A & 5,6 & 2,4 & 8 \\[3pt]
            B & 14 & 6 & 20 \\[3pt]
            C & 4,2 & 1,2 & 5,4 \\[3pt]
            D & 1,4 & 0,6 & 2 \\
            \bottomrule
        \end{tabular}
    \end{center}
    b) Calculamos la proporción \(m_{\ch{Fe}}/m_{\ch{O}}\): A \(\num{5.6}/\num{2.4} = \num{2.33}\); B \(14/6 = \num{2.33}\);
    C \(\num{4.2}/\num{1.2} = \num{3.5}\); D \(\num{1.4}/\num{0.6} = \num{2.33}\). Por la ley de Proust, A, B y D son el mismo óxido;
    C es otro diferente.

    c) En el óxido C, \(\dfrac{m_{\ch{Fe}}}{m_{\ch{O}}} = \num{3.5} \Rightarrow m_{\ch{O}} = \dfrac{21}{\num{3.5}} = \qty{6}{\gram}\) de oxígeno.
\end{example}

\begin{example}
    \textbf{Amoníaco.} Una muestra contiene \qty{85}{\gram} de amoníaco, \ch{NH3}. Calcula: a) la masa molecular; b) los moles;
    c) las moléculas; d) el número total de átomos; e) los átomos de nitrógeno y de hidrógeno; f) la composición centesimal.

    \textbf{Solución.} a) \(M(\ch{NH3}) = 14 + 3 \cdot 1 = \qty{17}{\uma}\) (\(M = \qty{17}{\gram\per\mole}\)).

    b) \(\qty{85}{\gram} \cdot \dfrac{\qty{1}{\mole}}{\qty{17}{\gram}} = \qty{5}{\mole}\).
    c) \(\qty{5}{\mole} \cdot \num{6.022e23} = \num{3.011e24}\) moléculas.

    d) Cada molécula tiene 4 átomos: \(4 \cdot \num{3.011e24} = \num{1.204e25}\) átomos.

    e) Un átomo de N y tres de H por molécula: \(\num{3.011e24}\) átomos de N y \(3 \cdot \num{3.011e24} = \num{9.033e24}\) de H.

    f) \(\%\,\ch{N} = \dfrac{14}{17} \cdot 100 = \qty{82.35}{\percent}\); \(\%\,\ch{H} = \dfrac{3}{17} \cdot 100 = \qty{17.65}{\percent}\).
\end{example}

\begin{example}
    \textbf{Leyes de los gases.} Un recipiente de \qty{5}{\litre} contiene un gas a \qty{2}{\atm} y \qty{60}{\degreeCelsius}.
    Indica qué ley aplicas en cada caso y calcula: a) el volumen si se duplica la temperatura absoluta a presión constante;
    b) la temperatura si la presión se reduce a la mitad sin cambiar el volumen; c) la presión si, a temperatura constante, el
    volumen pasa a ser \qty{25}{\litre}.

    \textbf{Solución.} Datos: \(V_1 = \qty{5}{\litre}\), \(p_1 = \qty{2}{\atm}\), \(T_1 = 60 + 273 = \qty{333}{\kelvin}\).

    a) Ley de Charles: \(\dfrac{V_1}{T_1} = \dfrac{V_2}{T_2} \Rightarrow V_2 = V_1 \dfrac{2 T_1}{T_1} = 2 \cdot 5 = \qty{10}{\litre}\).

    b) Ley de Gay-Lussac: \(\dfrac{p_1}{T_1} = \dfrac{p_2}{T_2} \Rightarrow T_2 = T_1 \dfrac{p_1/2}{p_1} = \dfrac{333}{2} = \qty{166.5}{\kelvin}\).

    c) Ley de Boyle-Mariotte: \(p_2 = \dfrac{p_1 V_1}{V_2} = \dfrac{2 \cdot 5}{25} = \qty{0.4}{\atm}\).
\end{example}

\begin{example}
    \textbf{¿Qué gas es?} a) La densidad de un gas es \qty{1.16}{\gram\per\litre} a \qty{750}{\mmHg} y \qty{15}{\degreeCelsius}.
    Averigua si es monóxido de carbono, \ch{CO}, monóxido de azufre, \ch{SO}, o amoníaco, \ch{NH3}. b) ¿Cuál sería la densidad de
    los otros dos gases en esas condiciones? c) ¿A qué presión la densidad del gas sería \qty{0.6}{\gram\per\litre} a
    \qty{30}{\degreeCelsius}?

    \textbf{Solución.} Datos: \(p = 750/760 = \qty{0.987}{\atm}\), \(T = \qty{288}{\kelvin}\).

    a) De \(p V = \dfrac{m}{M} R T\) y \(d = m/V\):
    \[ M = \frac{d R T}{p} = \frac{\num{1.16} \cdot \num{0.082} \cdot 288}{\num{0.987}} = \qty{27.8}{\gram\per\mole} \approx 28 \] Como \(M(\ch{CO}) = \qty{28}{\gram\per\mole}\), \(M(\ch{SO}) = 48\) y
    \(M(\ch{NH3}) = 17\), el gas es \textbf{monóxido de carbono}.

    b) \(d = \dfrac{p M}{R T}\): \ch{SO}: \(\dfrac{\num{0.987} \cdot 48}{\num{0.082} \cdot 288} = \qty{2.01}{\gram\per\litre}\);
    \ch{NH3}: \(\dfrac{\num{0.987} \cdot 17}{\num{0.082} \cdot 288} = \qty{0.71}{\gram\per\litre}\).

    c) \(p = \dfrac{d R T}{M} = \dfrac{\num{0.6} \cdot \num{0.082} \cdot 303}{28} = \qty{0.53}{\atm}\).
\end{example}

\begin{example}
    \textbf{Mezcla de vapores.} En un recipiente de \qty{10}{\litre} se introducen \qty{9.2}{\gram} de etanol, \ch{CH3CH2OH}, y
    \qty{11.6}{\gram} de acetona, \ch{CH3COCH3}. Después de cerrarlo, se calienta hasta \qty{280}{\degreeCelsius}, temperatura muy
    superior a los puntos de ebullición de ambos líquidos. Calcula: a) la presión parcial del etanol; b) la de la acetona; c) la
    presión total.

    \textbf{Solución.} Cada gas se comporta como si estuviera solo en los \qty{10}{\litre} (ley de Dalton). \(T = \qty{553}{\kelvin}\).
    \(n_{\text{etanol}} = \dfrac{\num{9.2}}{46} = \qty{0.2}{\mole}\); \(n_{\text{acetona}} = \dfrac{\num{11.6}}{58} = \qty{0.2}{\mole}\).

    a) y b) \(p = \dfrac{n R T}{V} = \dfrac{\num{0.2} \cdot \num{0.082} \cdot 553}{10} = \qty{0.907}{\atm}\) para cada uno, porque hay los
    mismos moles de los dos.

    c) \(p_T = \num{0.907} + \num{0.907} = \qty{1.81}{\atm}\). También se puede calcular con los moles totales:
    \(p_T = \dfrac{\num{0.4} \cdot \num{0.082} \cdot 553}{10} = \qty{1.81}{\atm}\).
\end{example}

%========================================================================================
\section{Ejercicios de aplicación}\label{sec:ud1-aplicacion}
%========================================================================================

\subsection{Leyes ponderales y volumétricas}

\begin{exercise}
    \(\star\) Completa la tabla de la reacción entre hierro y azufre para formar sulfuro de hierro(II),
    \(\ch{Fe(s) + S(s) -> FeS(s)}\), sabiendo que en la primera fila reaccionan exactamente las masas indicadas.
    \begin{center}
        \begin{tabular}{@{}ccc@{}}
            \toprule
            \(m\)\textbf{(Fe) (g)} & \(m\)\textbf{(S) (g)} & \(m\)\textbf{(FeS) (g)} \\
            \midrule
            55,8 & 32 & A \\[3pt]
            80 & B & C \\[3pt]
            D & 15 & E \\
            \bottomrule
        \end{tabular}
    \end{center}
    \begin{solu}
        \textcolor{blue}{Proporción \(m_{\ch{Fe}}/m_{\ch{S}} = \num{55.8}/32 = \num{1.744}\).}
        \(A = \num{55.8} + 32 = \textcolor{red}{\num{87.8}}\); \(B = 80/\num{1.744} = \textcolor{red}{\num{45.9}}\);
        \(C = 80 + \num{45.9} = \textcolor{red}{\num{125.9}}\); \(D = 15 \cdot \num{1.744} = \textcolor{red}{\num{26.2}}\);
        \(E = \num{26.2} + 15 = \textcolor{red}{\num{41.2}}\) (en gramos).
    \end{solu}
\end{exercise}

\begin{exercise}
    \(\star\star\) Al reaccionar completamente \qty{4}{\gram} de una sustancia A con \qty{6}{\gram} de otra B se forman
    \qty{10}{\gram} de un compuesto C. ¿Qué masa de C se obtendrá si se mezclan \qty{10}{\gram} de A con \qty{12}{\gram} de B?
    ¿Sobrará algún reactivo?
    \begin{solu}
        \textcolor{blue}{Proporción A:B = 4:6. Para 10 g de A harían falta 15 g de B, y solo hay 12 g: B se agota.} Con 12 g de B
        reaccionan \(12 \cdot 4/6 = \qty{8}{\gram}\) de A. Se forman \(8 + 12 = \textcolor{red}{\qty{20}{\gram}}\) de C y
        \textcolor{red}{sobran \qty{2}{\gram} de A}.
    \end{solu}
\end{exercise}

\begin{exercise}
    \(\star\star\) El azufre y el oxígeno forman dióxido de azufre combinándose en la proporción en masa 1:1. Completa la tabla
    (\(m_{\text{exc}}\) es la masa de reactivo que sobra) e indica qué ley ponderal se cumple.
    \begin{center}
        \begin{tabular}{@{}cccccc@{}}
            \toprule
            \textbf{Experiencia} & \(m\)\textbf{(S)} & \(m\)\textbf{(O)} & \(m\)\textbf{(\ch{SO2})} & \(m_{\text{exc}}\)\textbf{(S)} & \(m_{\text{exc}}\)\textbf{(O)} \\
            \midrule
            1 & 4 & 4 & 8 & 0 & 0 \\[3pt]
            2 & 5 & 7 & A & B & 2 \\[3pt]
            3 & C & 7 & 12 & 1 & D \\[3pt]
            4 & 9 & 4 & E & F & G \\
            \bottomrule
        \end{tabular}
    \end{center}
    \begin{solu}
        \textcolor{blue}{Ley de Proust: siempre reaccionan masas iguales de S y de O.}
        Exp. 2: reaccionan 5 g de S y 5 g de O: \(A = \textcolor{red}{10}\), \(B = \textcolor{red}{0}\).
        Exp. 3: 12 g de \ch{SO2} son 6 g de S y 6 g de O: \(D = 7 - 6 = \textcolor{red}{1}\), \(C = 6 + 1 = \textcolor{red}{7}\).
        Exp. 4: el O se agota: \(E = \textcolor{red}{8}\), \(F = 9 - 4 = \textcolor{red}{5}\), \(G = \textcolor{red}{0}\) (en gramos).
    \end{solu}
\end{exercise}

\begin{exercise}
    \(\star\star\) El plomo y el cloro forman dos compuestos. En una muestra del primero hay \qty{0.5836}{\gram} de Pb y
    \qty{0.200}{\gram} de Cl; en una del segundo, \qty{0.4377}{\gram} de Pb y \qty{0.300}{\gram} de Cl.
    a) Halla la proporción entre Pb y Cl en cada muestra. b) Si el primer compuesto es \ch{PbCl2}, ¿cuál es el segundo?
    c) ¿Cuántos gramos de cloro deben reaccionar con \qty{60}{\gram} de Pb para formar cada compuesto?
    \begin{solu}
        a) Primero: \(\num{0.5836}/\num{0.200} = \textcolor{red}{\num{2.918}}\); segundo: \(\num{0.4377}/\num{0.300} = \textcolor{red}{\num{1.459}}\).
        b) \textcolor{blue}{Con la misma masa de Pb, el segundo tiene el doble de cloro} (\(\num{2.918}/\num{1.459} = 2\)):
        \textcolor{red}{\ch{PbCl4}}.
        c) \(60/\num{2.918} = \textcolor{red}{\qty{20.6}{\gram}}\) de Cl para el primero y \(60/\num{1.459} = \textcolor{red}{\qty{41.1}{\gram}}\)
        para el segundo.
    \end{solu}
\end{exercise}

\begin{exercise}
    \(\star\star\) En el amoníaco, el nitrógeno y el hidrógeno están en la proporción en masa 14:3. a) ¿Qué masa de hidrógeno
    hay que añadir para que reaccionen \qty{120}{\gram} de nitrógeno? b) ¿Cuánto amoníaco se obtiene? c) ¿Cuánto amoníaco
    se obtiene con \qty{90}{\gram} de hidrógeno y \qty{40}{\gram} de nitrógeno? ¿Sobra algo?
    \begin{solu}
        a) \(m_{\ch{H}} = 120 \cdot 3/14 = \textcolor{red}{\qty{25.7}{\gram}}\). b) \(120 + \num{25.7} = \textcolor{red}{\qty{145.7}{\gram}}\) de \ch{NH3}.
        c) \textcolor{blue}{Para 40 g de N bastan \(40 \cdot 3/14 = \qty{8.57}{\gram}\) de H: el N se agota.} Se forman
        \(40 + \num{8.57} = \textcolor{red}{\qty{48.6}{\gram}}\) de \ch{NH3} y \textcolor{red}{sobran \qty{81.4}{\gram} de hidrógeno}.
    \end{solu}
\end{exercise}

\begin{exercise}
    \(\star\star\) El cobre y el oxígeno forman dos óxidos. En el óxido A la proporción en masa Cu:O es 3,97:1 y en el óxido B es
    7,94:1. a) Comprueba que se cumple la ley de las proporciones múltiples. b) Si A tiene un átomo de Cu por cada átomo de O,
    ¿cuál es la fórmula de B?
    \begin{solu}
        a) \textcolor{blue}{Con 1 g de O se combinan 3,97 g de Cu en A y 7,94 g en B:} \(\num{7.94}/\num{3.97} = \textcolor{red}{2}\), número entero
        sencillo. b) \textcolor{blue}{B tiene el doble de cobre por cada átomo de oxígeno:} \textcolor{red}{\ch{Cu2O}}.
    \end{solu}
\end{exercise}

\begin{exercise}
    \(\star\star\) Tras una excursión se recogen dos muestras de un mineral de cromo y oxígeno: la primera contiene
    \qty{52.0}{\gram} de Cr y \qty{24.0}{\gram} de O, y la segunda, \qty{39.0}{\gram} de Cr y \qty{18.0}{\gram} de O. ¿Se trata del
    mismo compuesto?
    \begin{solu}
        \(\num{52.0}/\num{24.0} = \num{2.167}\) y \(\num{39.0}/\num{18.0} = \num{2.167}\): \textcolor{red}{la proporción es la misma, luego es el
        mismo compuesto} (ley de Proust).
    \end{solu}
\end{exercise}

\begin{exercise}
    \(\star\) Un litro de hidrógeno reacciona con un litro de cloro para dar dos litros de cloruro de hidrógeno. ¿Cuántos litros
    de \ch{HCl} se obtienen si reaccionan \qty{6}{\litre} de hidrógeno con \qty{9}{\litre} de cloro? ¿Sobra algo? (Todos los
    volúmenes, en las mismas condiciones.)
    \begin{solu}
        \textcolor{blue}{Relación 1:1:2; el \ch{H2} se agota.} Se obtienen \(2 \cdot 6 = \textcolor{red}{\qty{12}{\litre}}\) de \ch{HCl} y
        \textcolor{red}{sobran \qty{3}{\litre} de \ch{Cl2}}.
    \end{solu}
\end{exercise}

\begin{exercise}
    \(\star\star\) Dos litros de nitrógeno reaccionan con seis litros de hidrógeno y dan cuatro litros de amoníaco; un litro de
    hidrógeno con uno de cloro da dos litros de cloruro de hidrógeno. ¿Contradicen estas reacciones la ley de Lavoisier?
    \begin{solu}
        \textcolor{red}{No.} \textcolor{blue}{La ley de Lavoisier se refiere a la masa, no al volumen. En la primera reacción el volumen
        disminuye (\(8 \to 4\) L) porque el número de moléculas disminuye (\(\ch{N2 + 3 H2 -> 2 NH3}\)), y en la segunda no cambia;
        pero en ambas los átomos se conservan y, por tanto, la masa también.}
    \end{solu}
\end{exercise}

\begin{exercise}
    \(\star\star\) El azufre gaseoso y el oxígeno forman tres óxidos; las relaciones de volumen S:\ch{O2} son 2:1, 1:1 y 2:3.
    a) ¿Qué volúmenes de azufre y de oxígeno hacen falta para obtener \qty{3}{\litre} de cada óxido (si se obtiene un volumen de
    óxido igual al de azufre)? b) ¿Cuál es la fórmula de cada óxido?
    \begin{solu}
        a) \textcolor{red}{3 L de S con 1,5 L, 3 L y 4,5 L de \ch{O2}}, respectivamente.
        b) \textcolor{blue}{Cada molécula de óxido lleva un átomo de S, y cada molécula de \ch{O2} aporta 2 átomos de O:}
        \textcolor{red}{\ch{SO}, \ch{SO2} y \ch{SO3}}.
    \end{solu}
\end{exercise}

\begin{exercise}
    \(\star\star\) El hidrógeno reacciona con el oxígeno para dar vapor de agua: \ch{2 H2(g) + O2(g) -> 2 H2O(g)}. Completa
    la tabla de volúmenes (en litros, en las mismas condiciones).
    \begin{center}
        \begin{tabular}{@{}ccccc@{}}
            \toprule
            \multicolumn{2}{c}{\textbf{Inicial}} & \multicolumn{3}{c}{\textbf{Final}} \\
            \cmidrule(r){1-2} \cmidrule(l){3-5}
            \ch{H2} & \ch{O2} & \ch{H2O} & \ch{H2} & \ch{O2} \\
            \midrule
            4 & 2 & 4 & 0 & 0 \\[3pt]
            6 & 4 & A & 0 & B \\[3pt]
            C & 3 & 6 & 0 & 0 \\[3pt]
            10 & 3 & D & E & 0 \\
            \bottomrule
        \end{tabular}
    \end{center}
    \begin{solu}
        \(A = \textcolor{red}{6}\), \(B = 4 - 3 = \textcolor{red}{1}\); \(C = \textcolor{red}{6}\); \textcolor{blue}{en la última fila se agota el \ch{O2}:}
        \(D = \textcolor{red}{6}\), \(E = 10 - 6 = \textcolor{red}{4}\).
    \end{solu}
\end{exercise}

\subsection{Cantidad de sustancia}

\begin{exercise}
    \(\star\) Indica si son verdaderas o falsas: a) La masa de un mol de \ch{O2} es \qty{16}{\gram}. b) En condiciones normales,
    un mol de \ch{CO2} y un mol de \ch{CH4} ocupan el mismo volumen. c) Un mol de \ch{H2O} y un mol de \ch{NH3} contienen el mismo
    número de átomos.
    \begin{solu}
        a) \textcolor{red}{F}: es \qty{32}{\gram}. b) \textcolor{red}{V}: ley de Avogadro. c) \textcolor{red}{F}: el mismo número de moléculas,
        pero cada molécula de agua tiene 3 átomos y cada una de amoníaco, 4.
    \end{solu}
\end{exercise}

\begin{exercise}
    \(\star\) Completa la tabla.
    \begin{center}
        \begin{tabular}{@{}cccc@{}}
            \toprule
            \textbf{Sustancia} & \textbf{Moles} & \textbf{Masa (g)} & \textbf{Moléculas} \\
            \midrule
            \ch{H2} & 3 & A & B \\[3pt]
            \ch{N2} & C & 70 & D \\[3pt]
            \ch{O2} & E & F & \num{2.4e24} \\[3pt]
            \ch{H2O} & 1,5 & G & H \\
            \bottomrule
        \end{tabular}
    \end{center}
    \begin{solu}
        \(A = 3 \cdot 2 = \textcolor{red}{6}\); \(B = 3 \cdot \num{6.022e23} = \textcolor{red}{\num{1.81e24}}\);
        \(C = 70/28 = \textcolor{red}{\num{2.5}}\); \(D = \textcolor{red}{\num{1.51e24}}\);
        \(E = \num{2.4e24}/\num{6.022e23} = \textcolor{red}{\num{3.99}}\); \(F = \num{3.99} \cdot 32 = \textcolor{red}{\num{127.5}}\);
        \(G = \num{1.5} \cdot 18 = \textcolor{red}{27}\); \(H = \textcolor{red}{\num{9.03e23}}\).
    \end{solu}
\end{exercise}

\begin{exercise}
    \(\star\) La cafeína, \ch{C8H10N4O2}, está presente en el café, el té y muchos refrescos. Un frasco de un laboratorio
    contiene \qty{50}{\gram} de cafeína. Calcula: a) los moles; b) las moléculas; c) la composición centesimal; d) los átomos de
    carbono y de nitrógeno.
    \begin{solu}
        \(M = \qty{194}{\gram\per\mole}\). a) \(50/194 = \textcolor{red}{\qty{0.258}{\mole}}\). b) \(\textcolor{red}{\num{1.55e23}}\) moléculas.
        c) \textcolor{red}{C \qty{49.5}{\percent}, H \qty{5.15}{\percent}, N \qty{28.9}{\percent}, O \qty{16.5}{\percent}}.
        d) C: \(8 \cdot \num{1.55e23} = \textcolor{red}{\num{1.24e24}}\); N: \(4 \cdot \num{1.55e23} = \textcolor{red}{\num{6.21e23}}\).
    \end{solu}
\end{exercise}

\begin{exercise}
    \(\star\star\) La nicotina, \ch{C10H14N2}, es la sustancia adictiva del tabaco. Una muestra contiene \num{2.5e25} átomos de
    carbono. Calcula: a) las moléculas de nicotina; b) la composición centesimal; c) los átomos de nitrógeno; d) los moles;
    e) la masa de la muestra.
    \begin{solu}
        a) \(\num{2.5e25}/10 = \textcolor{red}{\num{2.5e24}}\). b) \(M = \qty{162}{\gram\per\mole}\): \textcolor{red}{C \qty{74.1}{\percent}, H \qty{8.6}{\percent},
        N \qty{17.3}{\percent}}. c) \(2 \cdot \num{2.5e24} = \textcolor{red}{\num{5e24}}\). d) \(\num{2.5e24}/\num{6.022e23} = \textcolor{red}{\qty{4.15}{\mole}}\).
        e) \(\num{4.15} \cdot 162 = \textcolor{red}{\qty{673}{\gram}}\).
    \end{solu}
\end{exercise}

\begin{exercise}
    \(\star\star\) ¿Cuál de estas muestras tiene mayor número de átomos? a) \qty{25}{\gram} de Fe. b) \num{4e23} moléculas de
    \ch{NH3}. c) \qty{1.2}{\mole} de \ch{CO2}.
    \begin{solu}
        a) \(25/\num{55.8} \cdot \num{6.022e23} = \num{2.70e23}\); b) \(4 \cdot \num{4e23} = \num{1.6e24}\);
        c) \(3 \cdot \num{1.2} \cdot \num{6.022e23} = \num{2.17e24}\). \textcolor{red}{La muestra c}.
    \end{solu}
\end{exercise}

\begin{exercise}
    \(\star\star\) Una gota de agua ocupa unos \qty{0.05}{\cubic\centi\metre}. ¿Cuántas moléculas de agua hay en 15 gotas? (Densidad
    del agua: \qty{1}{\gram\per\cubic\centi\metre}.)
    \begin{solu}
        \(m = 15 \cdot \num{0.05} = \qty{0.75}{\gram}\); \(N = \dfrac{\num{0.75}}{18} \cdot \num{6.022e23} = \textcolor{red}{\num{2.51e22}}\) moléculas.
    \end{solu}
\end{exercise}

\begin{exercise}
    \(\star\star\) Calcula: a) los gramos de fósforo que hay en \qty{2.4}{\mole} de fosfano, \ch{PH3}; b) los átomos de oxígeno en
    \qty{9.45}{\gram} de ácido nítrico, \ch{HNO3}; c) los moles de átomos de oxígeno en \qty{3.5}{\mole} de \ch{Al2(SO4)3}.
    \begin{solu}
        a) \(\num{2.4} \cdot 31 = \textcolor{red}{\qty{74.4}{\gram}}\). b) \(M = \qty{63}{\gram\per\mole}\): \(\dfrac{\num{9.45}}{63} \cdot 3 \cdot \num{6.022e23} =
        \textcolor{red}{\num{2.71e23}}\). c) \textcolor{blue}{12 átomos de O por fórmula:} \(12 \cdot \num{3.5} = \textcolor{red}{\qty{42}{\mole}}\).
    \end{solu}
\end{exercise}

\subsection{Fórmula empírica y molecular}

\begin{exercise}
    \(\star\) Calcula la composición centesimal del fosfato de calcio, \ch{Ca3(PO4)2}. (Ca = 40; P = 31.)
    \begin{solu}
        \(M = 3 \cdot 40 + 2 \cdot 31 + 8 \cdot 16 = \qty{310}{\gram\per\mole}\).
        \textcolor{red}{Ca \qty{38.7}{\percent}; P \qty{20.0}{\percent}; O \qty{41.3}{\percent}}.
    \end{solu}
\end{exercise}

\begin{exercise}
    \(\star\star\) Una roca de \qty{40}{\kilogram} contiene un 80\,\% de carbonato de calcio, \ch{CaCO3}. Calcula: a) la masa de
    \ch{CaCO3}; b) la masa de oxígeno de la roca, si todo procede del carbonato.
    \begin{solu}
        a) \(\num{0.80} \cdot 40 = \textcolor{red}{\qty{32}{\kilogram}}\). b) \(M(\ch{CaCO3}) = \qty{100}{\gram\per\mole}\), con 48 g de O:
        \(32 \cdot \dfrac{48}{100} = \textcolor{red}{\qty{15.4}{\kilogram}}\).
    \end{solu}
\end{exercise}

\begin{exercise}
    \(\star\star\) Las plantas necesitan magnesio para fabricar la clorofila \textit{a}, de fórmula \ch{C55H72MgN4O5}. Calcula el
    porcentaje en masa de magnesio. (Mg = 24,3.)
    \begin{solu}
        \(M = \qty{892.3}{\gram\per\mole}\); \(\%\,\ch{Mg} = \dfrac{\num{24.3}}{\num{892.3}} \cdot 100 = \textcolor{red}{\qty{2.72}{\percent}}\).
    \end{solu}
\end{exercise}

\begin{exercise}
    \(\star\) Un óxido de hierro tiene un 72,4\,\% de hierro y un 27,6\,\% de oxígeno. Calcula su fórmula empírica.
    \begin{solu}
        Fe \(\num{72.4}/\num{55.8} = \num{1.297}\); O \(\num{27.6}/16 = \num{1.725}\) \(\rightarrow\) 1 : 1,33 \(\rightarrow\) (×3) 3 : 4.
        \textcolor{red}{\ch{Fe3O4}} (magnetita).
    \end{solu}
\end{exercise}

\begin{exercise}
    \(\star\) La glicina, el aminoácido más sencillo, contiene un 32,0\,\% de C, un 18,7\,\% de N, un 42,6\,\% de O y el resto de
    H. Calcula su fórmula empírica.
    \begin{solu}
        H: \qty{6.7}{\percent}. C \(\num{2.67}\); H \(\num{6.7}\); N \(\num{1.34}\); O \(\num{2.66}\) \(\rightarrow\) 2 : 5 : 1 : 2.
        \textcolor{red}{\ch{C2H5NO2}}.
    \end{solu}
\end{exercise}

\begin{exercise}
    \(\star\star\) La cafeína tiene un 49,5\,\% de C, un 5,15\,\% de H, un 28,9\,\% de N y un 16,5\,\% de O. Su masa molecular es 194.
    Determina su fórmula empírica y su fórmula molecular.
    \begin{solu}
        C \(\num{4.13}\); H \(\num{5.15}\); N \(\num{2.06}\); O \(\num{1.03}\) \(\rightarrow\) 4 : 5 : 2 : 1: empírica \textcolor{red}{\ch{C4H5N2O}}
        (\(M = 97\)). \(\alpha = 194/97 = 2\): molecular \textcolor{red}{\ch{C8H10N4O2}}.
    \end{solu}
\end{exercise}

\begin{exercise}
    \(\star\) Un compuesto gaseoso de carbono, hidrógeno y oxígeno tiene de fórmula molecular \ch{C2H6O}. Calcula: a) su fórmula
    empírica; b) los moles de moléculas en \qty{23}{\gram}; c) los gramos de carbono en \qty{23}{\gram}; d) los moles de átomos de
    oxígeno en \qty{23}{\gram}; e) los átomos de hidrógeno en \qty{23}{\gram}.
    \begin{solu}
        a) \textcolor{red}{\ch{C2H6O}} (no se puede simplificar). \(M = \qty{46}{\gram\per\mole}\). b) \(23/46 = \textcolor{red}{\qty{0.5}{\mole}}\).
        c) \(23 \cdot 24/46 = \textcolor{red}{\qty{12}{\gram}}\). d) \textcolor{red}{\qty{0.5}{\mole}}. e) \(6 \cdot \num{0.5} \cdot \num{6.022e23} = \textcolor{red}{\num{1.81e24}}\).
    \end{solu}
\end{exercise}

\begin{exercise}
    \(\star\star\) Un compuesto de nitrógeno y oxígeno contiene \qty{1.40}{\gram} de N por cada \qty{4.00}{\gram} de O. Su masa
    molecular es 108. Determina su fórmula molecular.
    \begin{solu}
        N \(\num{1.40}/14 = \num{0.10}\); O \(\num{4.00}/16 = \num{0.25}\) \(\rightarrow\) 1 : 2,5 \(\rightarrow\) 2 : 5: empírica \ch{N2O5}
        (\(M = 108\)). \(\alpha = 1\): \textcolor{red}{\ch{N2O5}}.
    \end{solu}
\end{exercise}

\begin{exercise}
    \(\star\star\) El benzoato de sodio (E-211) es un conservante alimentario. Contiene un 58,33\,\% de C, un 3,47\,\% de H, un
    22,22\,\% de O y un 15,97\,\% de Na. Su masa molecular relativa es 144. ¿Cuál es su fórmula molecular?
    \begin{solu}
        C \(\num{4.86}\); H \(\num{3.47}\); O \(\num{1.389}\); Na \(\num{0.694}\) \(\rightarrow\) 7 : 5 : 2 : 1; \(M = 144\), \(\alpha = 1\):
        \textcolor{red}{\ch{C7H5O2Na}}.
    \end{solu}
\end{exercise}

\begin{exercise}
    \(\star\star\) El paracetamol tiene un 63,58\,\% de C, un 5,96\,\% de H, un 9,27\,\% de N y un 21,19\,\% de O. Su masa molar es
    \qty{151}{\gram\per\mole}. ¿Cuál es su fórmula empírica? ¿Y la molecular?
    \begin{solu}
        C \(\num{5.30}\); H \(\num{5.96}\); N \(\num{0.662}\); O \(\num{1.324}\) \(\rightarrow\) 8 : 9 : 1 : 2.
        \textcolor{red}{Empírica y molecular: \ch{C8H9NO2}} (\(M = 151\), \(\alpha = 1\)).
    \end{solu}
\end{exercise}

\subsection{Leyes de los gases}

\begin{exercise}
    \(\star\) ¿Cómo debe variar la presión de un gas para que su volumen se reduzca a la tercera parte, a temperatura constante?
    \begin{solu}
        \textcolor{blue}{Boyle: \(p\) y \(V\) son inversamente proporcionales.} \textcolor{red}{La presión debe triplicarse.}
    \end{solu}
\end{exercise}

\begin{exercise}
    \(\star\) Un gas está en un recipiente rígido. ¿Qué le pasa a su presión si la temperatura absoluta se triplica?
    \begin{solu}
        \textcolor{blue}{Gay-Lussac, a volumen constante \(p/T\) es constante:} \textcolor{red}{la presión se triplica}.
    \end{solu}
\end{exercise}

\begin{exercise}
    \(\star\) Justifica si son verdaderas o falsas: a) Al duplicar la temperatura absoluta de un gas sin variar el volumen, la
    presión se duplica. b) Al calentar un gas a presión constante, su volumen disminuye. c) Si se reduce a la mitad el volumen
    de un gas a presión constante, su temperatura absoluta se duplica.
    \begin{solu}
        a) \textcolor{red}{V} (Gay-Lussac). b) \textcolor{red}{F}: aumenta (Charles). c) \textcolor{red}{F}: se reduce a la mitad (Charles).
    \end{solu}
\end{exercise}

\begin{exercise}
    \(\star\star\) Un globo aerostático, sin variar su volumen, reduce su presión a la mitad durante el ascenso. ¿Cómo ha
    cambiado su temperatura?
    \begin{solu}
        \(\dfrac{p_1}{T_1} = \dfrac{p_1/2}{T_2} \Rightarrow \textcolor{red}{T_2 = T_1/2}\): \textcolor{red}{la temperatura absoluta se reduce a la mitad.}
    \end{solu}
\end{exercise}

\begin{exercise}
    \(\star\star\) \qty{5}{\litre} de \ch{NH3}(g) están a \qty{4}{\atm}. Si, a temperatura constante, el volumen se multiplica por 4,
    ¿cuál será la presión?
    \begin{solu}
        \(p_2 = \dfrac{4 \cdot 5}{20} = \textcolor{red}{\qty{1}{\atm}}\).
    \end{solu}
\end{exercise}

\begin{exercise}
    \(\star\star\) Un globo inflado a \qty{2}{\atm} ocupa \qty{15}{\litre} y explota si su volumen supera los \qty{25}{\cubic\deci\metre}.
    Calcula su volumen si la presión baja hasta \qty{1.5}{\atm} a temperatura constante. ¿Explotará?
    \begin{solu}
        \(V_2 = \dfrac{2 \cdot 15}{\num{1.5}} = \textcolor{red}{\qty{20}{\litre}}\); \textcolor{red}{no explota} (\(20 < 25\)).
    \end{solu}
\end{exercise}

\begin{exercise}
    \(\star\star\) Un recipiente de \qty{250}{\cubic\centi\metre} contiene un gas a \qty{25}{\degreeCelsius} y \qty{3000}{\mmHg}.
    Manteniendo el volumen, ¿cuál será la presión si la temperatura baja hasta \qty{-10}{\degreeCelsius}?
    \begin{solu}
        \(p_2 = 3000 \cdot \dfrac{263}{298} = \qty{2648}{\mmHg} = \textcolor{red}{\qty{3.48}{\atm}}\).
    \end{solu}
\end{exercise}

\begin{exercise}
    \(\star\star\) Un contenedor de oxígeno resiste hasta \qty{40}{\atm}. Se llena a \qty{20}{\degreeCelsius} con una presión de
    \qty{11400}{\mmHg}. a) Representa la presión frente a la temperatura. b) ¿Qué ley se refleja en la gráfica? c) Calcula la
    presión si la temperatura sube a \qty{45}{\degreeCelsius}. d) ¿A qué temperatura explotaría?
    \begin{solu}
        \(p_1 = 11400/760 = \qty{15}{\atm}\), \(T_1 = \qty{293}{\kelvin}\).
        a) \textcolor{blue}{Una recta que pasa por el origen (con \(T\) en kelvin) de pendiente \(15/293 = \qty{0.0512}{\atm\per\kelvin}\):}
        \begin{center}
            \begin{tikzpicture}
                \begin{axis}[grafica, xlabel={\(T\) (\unit{\kelvin})}, ylabel={\(p\) (\unit{\atm})}, xmin=0, xmax=850, ymin=0, ymax=45]
                    \addplot[thick, blue, domain=0:800] {15/293*x};
                    \addplot[dashed, red] coordinates {(0,40) (781,40) (781,0)};
                    \addplot[only marks, mark=*, mark size=1.5pt] coordinates {(293,15) (318,16.28)};
                \end{axis}
            \end{tikzpicture}
        \end{center}
        b) \textcolor{red}{Ley de Gay-Lussac.} c) \(p_2 = 15 \cdot \dfrac{318}{293} = \textcolor{red}{\qty{16.3}{\atm}}\).
        d) \(T = 293 \cdot \dfrac{40}{15} = \qty{781}{\kelvin} = \textcolor{red}{\qty{508}{\degreeCelsius}}\).
    \end{solu}
\end{exercise}

\begin{exercise}
    \(\star\star\) La relación entre el volumen y la temperatura de un gas a presión constante es \(V/T = \qty{0.5}{\litre\per\kelvin}\).
    a) Haz una tabla \(V\)-\(T\) con cinco temperaturas. b) Represéntala. c) ¿Qué relación hay entre \(V\) y \(T\)? ¿Qué ley la describe?
    \begin{solu}
        a) Por ejemplo: \(T = 100, 200, 300, 400, 500\) K \(\rightarrow\) \(V = 50, 100, 150, 200, 250\) L.
        b) \textcolor{blue}{Recta que pasa por el origen de pendiente \qty{0.5}{\litre\per\kelvin}.}
        c) \textcolor{red}{Son directamente proporcionales: ley de Charles.}
    \end{solu}
\end{exercise}

\begin{exercise}
    \(\star\star\) Un globo aerostático tiene un volumen de \qty{2.5e5}{\cubic\deci\metre} antes de despegar, a \qty{30}{\degreeCelsius}.
    La temperatura baja \qty{1}{\degreeCelsius} cada \qty{200}{\metre} de altura. Calcula su volumen a \qty{2}{\kilo\metre} de altura,
    suponiendo que la presión no cambia.
    \begin{solu}
        \textcolor{blue}{A 2 km la temperatura ha bajado \qty{10}{\degreeCelsius}: \(T_2 = \qty{293}{\kelvin}\).}
        \(V_2 = \num{2.5e5} \cdot \dfrac{293}{303} = \textcolor{red}{\qty{2.42e5}{\litre}}\).
    \end{solu}
\end{exercise}

\subsection{Teoría cinético-molecular}

\begin{exercise}
    \(\star\) Explica con la teoría cinético-molecular cómo cambian las condiciones de las partículas de un gas encerrado si:
    a) reducimos el volumen del recipiente; b) aumentamos la temperatura a volumen constante; c) añadimos al recipiente un gas
    inerte (que no reacciona).
    \begin{solu}
        a) \textcolor{red}{Las partículas tienen menos espacio y chocan más a menudo con las paredes: sube la presión} (su velocidad media
        no cambia si \(T\) es constante). b) \textcolor{red}{Se mueven más deprisa y chocan con más frecuencia y más fuerza: sube la
        presión.} c) \textcolor{red}{Hay más partículas chocando con las paredes: aumenta la presión total}, aunque la presión parcial
        del primer gas no cambia.
    \end{solu}
\end{exercise}

\begin{exercise}
    \(\star\star\) Lucía se encuentra un globo y decide inflarlo más, pero acaba explotando. Argumenta cuál de estas afirmaciones
    explica lo que ha pasado: a) Al aumentar la cantidad de aire aumenta la temperatura del globo. b) El volumen del globo
    aumenta porque disminuye la presión. c) Al aumentar la cantidad de aire aumenta la presión interior.
    \begin{solu}
        a) \textcolor{red}{F}. b) \textcolor{red}{F}. c) \textcolor{red}{V}: \textcolor{blue}{con más partículas hay más choques contra la pared del globo; la goma se
        estira hasta que no resiste la presión.}
    \end{solu}
\end{exercise}

\subsection{Ecuación de los gases ideales}

\begin{exercise}
    \(\star\) ¿Qué volumen ocupan \qty{40}{\gram} de oxígeno en condiciones normales? ¿Y a \qty{1}{\atm} y \qty{30}{\degreeCelsius}?
    \begin{solu}
        \(n = 40/32 = \qty{1.25}{\mole}\). En c.\,n.: \(\num{1.25} \cdot \num{22.4} = \textcolor{red}{\qty{28}{\litre}}\).
        A \qty{30}{\degreeCelsius}: \(\num{1.25} \cdot \num{0.082} \cdot 303 = \textcolor{red}{\qty{31.1}{\litre}}\).
    \end{solu}
\end{exercise}

\begin{exercise}
    \(\star\star\) Determina el volumen, en unidades del SI, de un recipiente que contiene \qty{2.5}{\mole} de un gas a
    \qty{800}{\mmHg} y \qty{60}{\degreeCelsius}.
    \begin{solu}
        \(V = \dfrac{\num{2.5} \cdot \num{0.082} \cdot 333}{800/760} = \qty{64.9}{\litre} = \textcolor{red}{\qty{6.49e-2}{\cubic\metre}}\).
    \end{solu}
\end{exercise}

\begin{exercise}
    \(\star\star\) Halla la presión en un recipiente de \qty{30}{\cubic\metre} en el que hay \qty{56}{\gram} de nitrógeno, \ch{N2}, a
    \qty{-20}{\degreeCelsius}.
    \begin{solu}
        \(n = \qty{2}{\mole}\); \(V = \qty{30000}{\litre}\); \(p = \dfrac{2 \cdot \num{0.082} \cdot 253}{30000} = \textcolor{red}{\qty{1.38e-3}{\atm}}\):
        \textcolor{blue}{un gas muy enrarecido}.
    \end{solu}
\end{exercise}

\begin{exercise}
    \(\star\star\) Tenemos \qty{3}{\litre} de dióxido de carbono a \qty{820}{\mmHg} y \qty{27}{\degreeCelsius}. Calcula: a) el volumen que
    ocuparía en condiciones normales; b) los moles y las moléculas.
    \begin{solu}
        b) \(n = \dfrac{(820/760) \cdot 3}{\num{0.082} \cdot 300} = \textcolor{red}{\qty{0.132}{\mole}}\); \(N = \textcolor{red}{\num{7.92e22}}\) moléculas.
        a) \(V = \num{0.132} \cdot \num{22.4} = \textcolor{red}{\qty{2.95}{\litre}}\).
    \end{solu}
\end{exercise}

\begin{exercise}
    \(\star\star\) Una bombona de \qty{26}{\litre} contiene butano, \ch{C4H10}, a \qty{20}{\degreeCelsius}. Llena tiene una masa de
    \qty{24.0}{\kilogram}, y vacía, de \qty{22.5}{\kilogram}. a) Calcula la masa molar del butano y la masa de butano de la bombona.
    b) ¿Qué presión habría dentro si todo el butano estuviera en estado gaseoso?
    \begin{solu}
        a) \(M = 4 \cdot 12 + 10 = \textcolor{red}{\qty{58}{\gram\per\mole}}\); \(m = \textcolor{red}{\qty{1.5}{\kilogram}}\).
        b) \(n = 1500/58 = \qty{25.9}{\mole}\); \(p = \dfrac{\num{25.9} \cdot \num{0.082} \cdot 293}{26} = \textcolor{red}{\qty{23.9}{\atm}}\).
        \textcolor{blue}{(En realidad, a esa presión el butano se licúa: por eso las bombonas contienen butano líquido.)}
    \end{solu}
\end{exercise}

\begin{exercise}
    \(\star\) Un recipiente de \qty{3}{\cubic\deci\metre} contiene \qty{8}{\gram} de nitrógeno a \qty{30}{\degreeCelsius}. Calcula la presión.
    \begin{solu}
        \(p = \dfrac{(8/28) \cdot \num{0.082} \cdot 303}{3} = \textcolor{red}{\qty{2.37}{\atm}}\).
    \end{solu}
\end{exercise}

\begin{exercise}
    \(\star\) ¿Cuántos moles de \ch{CO2} hay en \qty{66}{\gram}? ¿Qué volumen ocupan a \qty{1}{\atm} y \qty{0}{\degreeCelsius}?
    \begin{solu}
        \(n = 66/44 = \textcolor{red}{\qty{1.5}{\mole}}\); \(V = \num{1.5} \cdot \num{22.4} = \textcolor{red}{\qty{33.6}{\litre}}\).
    \end{solu}
\end{exercise}

\begin{exercise}
    \(\star\star\) El kriptón, Kr(g), se almacena en una botella de \qty{15}{\litre} a \qty{10}{\atm} y \qty{20}{\degreeCelsius}. Calcula la
    masa de kriptón. (Kr = 83,8.)
    \begin{solu}
        \(n = \dfrac{10 \cdot 15}{\num{0.082} \cdot 293} = \qty{6.24}{\mole}\); \(m = \num{6.24} \cdot \num{83.8} = \textcolor{red}{\qty{523}{\gram}}\).
    \end{solu}
\end{exercise}

\begin{exercise}
    \(\star\star\) La presión dentro de un bote de laca es de \qty{4}{\bar} a \qty{25}{\degreeCelsius}. ¿Cuántos moles de gas hay en un bote de
    \qty{400}{\milli\litre}? ¿Y en uno de \qty{250}{\cubic\centi\metre}?
    \begin{solu}
        \(p = 4/\num{1.013} = \qty{3.95}{\atm}\). \(n_1 = \dfrac{\num{3.95} \cdot \num{0.4}}{\num{0.082} \cdot 298} = \textcolor{red}{\qty{0.065}{\mole}}\);
        \(n_2 = \textcolor{red}{\qty{0.040}{\mole}}\).
    \end{solu}
\end{exercise}

\begin{exercise}
    \(\star\star\) Una cantidad de gas ocupa \qty{1.5}{\litre} a cierta temperatura. ¿Cómo podríamos guardar esa misma cantidad en un
    recipiente de \qty{1}{\litre}? ¿Y en uno de \qty{5}{\litre}? La temperatura no cambia.
    \begin{solu}
        \textcolor{red}{En el de 1 L, comprimiéndolo: su presión aumenta 1,5 veces. En el de 5 L, el gas se expande hasta ocupar todo el
        recipiente y su presión se reduce a \(1{,}5/5 = 0{,}3\) veces la inicial} (ley de Boyle).
    \end{solu}
\end{exercise}

\subsection{Densidad de los gases}

\begin{exercise}
    \(\star\) ¿Cuál es la densidad del dióxido de azufre, \ch{SO2}, a \qty{120}{\degreeCelsius} y \qty{1.5}{\atm}?
    \begin{solu}
        \(d = \dfrac{p M}{R T} = \dfrac{\num{1.5} \cdot 64}{\num{0.082} \cdot 393} = \textcolor{red}{\qty{2.98}{\gram\per\litre}}\).
    \end{solu}
\end{exercise}

\begin{exercise}
    \(\star\) Calcula la densidad del metano, \ch{CH4}: a) a \qty{80}{\degreeCelsius} y \qty{900}{\mmHg}; b) en condiciones normales.
    \begin{solu}
        a) \(d = \dfrac{(900/760) \cdot 16}{\num{0.082} \cdot 353} = \textcolor{red}{\qty{0.655}{\gram\per\litre}}\).
        b) \(d = 16/\num{22.4} = \textcolor{red}{\qty{0.714}{\gram\per\litre}}\).
    \end{solu}
\end{exercise}

\begin{exercise}
    \(\star\star\) La densidad de un gas es \qty{2.5}{\gram\per\litre} a \qty{20}{\degreeCelsius} y \qty{1}{\atm}. ¿Cuál será a
    \qty{60}{\degreeCelsius} y \qty{1.5}{\atm}?
    \begin{solu}
        \textcolor{blue}{\(d\) es proporcional a \(p/T\):} \(d_2 = \num{2.5} \cdot \dfrac{\num{1.5}}{1} \cdot \dfrac{293}{333} = \textcolor{red}{\qty{3.30}{\gram\per\litre}}\).
    \end{solu}
\end{exercise}

\begin{exercise}
    \(\star\) Una muestra de \qty{8}{\gram} de un gas ocupa \qty{3}{\litre} a \qty{1.8}{\atm} y \qty{27}{\degreeCelsius}. Calcula su masa molar.
    \begin{solu}
        \(M = \dfrac{m R T}{p V} = \dfrac{8 \cdot \num{0.082} \cdot 300}{\num{1.8} \cdot 3} = \textcolor{red}{\qty{36.4}{\gram\per\mole}}\).
    \end{solu}
\end{exercise}

\begin{exercise}
    \(\star\star\) Para un experimento se necesita nitrógeno con una densidad de \qty{2.2}{\gram\per\litre} a \qty{25}{\degreeCelsius}. ¿A qué
    presión debe estar, en \unit{\mmHg} y en \unit{\pascal}?
    \begin{solu}
        \(p = \dfrac{d R T}{M} = \dfrac{\num{2.2} \cdot \num{0.082} \cdot 298}{28} = \qty{1.92}{\atm} = \textcolor{red}{\qty{1459}{\mmHg}} =
        \textcolor{red}{\qty{1.95e5}{\pascal}}\).
    \end{solu}
\end{exercise}

\begin{exercise}
    \(\star\star\) La densidad de un gas en condiciones normales es \qty{0.76}{\gram\per\litre}. ¿Cuál es el gas? a) Monóxido de carbono.
    b) Dióxido de carbono. c) Amoníaco.
    \begin{solu}
        \(M = \num{0.76} \cdot \num{22.4} = \qty{17}{\gram\per\mole}\): \textcolor{red}{c) amoníaco}.
    \end{solu}
\end{exercise}

\begin{exercise}
    \(\star\star\) \qty{80}{\litre} de un gas tienen una masa de \qty{150}{\gram} a \qty{1e5}{\pascal} y \qty{273}{\kelvin}. Calcula: a) la densidad en
    esas condiciones; b) la densidad si, manteniendo la presión, la temperatura sube a \qty{80}{\degreeCelsius}.
    \begin{solu}
        a) \(d = 150/80 = \textcolor{red}{\qty{1.88}{\gram\per\litre}}\). b) \(d_2 = \num{1.875} \cdot \dfrac{273}{353} = \textcolor{red}{\qty{1.45}{\gram\per\litre}}\).
    \end{solu}
\end{exercise}

\subsection{Mezclas de gases y presiones parciales}

\begin{exercise}
    \(\star\) En un recipiente se mezclan \qty{35}{\gram} de nitrógeno y \qty{48}{\gram} de oxígeno. La presión total es de \qty{2}{\atm}.
    Calcula la presión parcial de cada gas.
    \begin{solu}
        \(n_{\ch{N2}} = \qty{1.25}{\mole}\), \(n_{\ch{O2}} = \qty{1.5}{\mole}\), \(n_T = \qty{2.75}{\mole}\).
        \(p_{\ch{N2}} = \dfrac{\num{1.25}}{\num{2.75}} \cdot 2 = \textcolor{red}{\qty{0.909}{\atm}}\); \(p_{\ch{O2}} = \textcolor{red}{\qty{1.091}{\atm}}\).
    \end{solu}
\end{exercise}

\begin{exercise}
    \(\star\star\) Se introducen \qty{6}{\gram} de hidrógeno y \qty{20}{\gram} de helio en un recipiente de \qty{25}{\litre} a
    \qty{20}{\degreeCelsius}. Calcula: a) la presión parcial de cada gas; b) la presión total.
    \begin{solu}
        a) \(p_{\ch{H2}} = \dfrac{3 \cdot \num{0.082} \cdot 293}{25} = \textcolor{red}{\qty{2.88}{\atm}}\); \(p_{\ch{He}} = \dfrac{5 \cdot \num{0.082} \cdot 293}{25} =
        \textcolor{red}{\qty{4.81}{\atm}}\). b) \(\textcolor{red}{\qty{7.69}{\atm}}\).
    \end{solu}
\end{exercise}

\begin{exercise}
    \(\star\star\) Una atmósfera inerte está formada por \qty{176}{\gram} de \ch{CO2}, \qty{84}{\gram} de \ch{N2} y \qty{26.3}{\gram} de xenón. Si
    la presión es de \qty{2280}{\mmHg}, calcula la presión parcial de cada componente. (Xe = 131,3.)
    \begin{solu}
        \(n\): \ch{CO2} 4; \ch{N2} 3; Xe \num{0.200}; \(n_T = \num{7.20}\). \(p_T = 2280/760 = \qty{3}{\atm}\).
        \(p_{\ch{CO2}} = \textcolor{red}{\qty{1.67}{\atm}}\); \(p_{\ch{N2}} = \textcolor{red}{\qty{1.25}{\atm}}\); \(p_{\ch{Xe}} = \textcolor{red}{\qty{0.083}{\atm}}\).
    \end{solu}
\end{exercise}

\begin{exercise}
    \(\star\star\) En un recipiente de \qty{30}{\litre} a \qty{0}{\degreeCelsius} hay \qty{66}{\gram} de dióxido de carbono. Se añaden
    \qty{0.5}{\mole} de monóxido de carbono. Calcula: a) la presión final de la mezcla; b) la presión parcial de cada componente.
    \begin{solu}
        \(n_{\ch{CO2}} = \qty{1.5}{\mole}\); \(n_T = \qty{2}{\mole}\).
        a) \(p_T = \dfrac{2 \cdot \num{0.082} \cdot 273}{30} = \textcolor{red}{\qty{1.49}{\atm}}\).
        b) \(p_{\ch{CO2}} = \textcolor{red}{\qty{1.12}{\atm}}\); \(p_{\ch{CO}} = \textcolor{red}{\qty{0.37}{\atm}}\).
    \end{solu}
\end{exercise}

\begin{exercise}
    \(\star\) La presión de una mezcla de dos gases A y B es de \qty{1800}{\mmHg} a \qty{20}{\degreeCelsius}. Si la fracción molar de A es
    0,35, calcula las presiones parciales.
    \begin{solu}
        \(p_T = \qty{2.37}{\atm}\). \(p_A = \num{0.35} \cdot \num{2.37} = \textcolor{red}{\qty{0.83}{\atm}}\); \(p_B = \textcolor{red}{\qty{1.54}{\atm}}\).
    \end{solu}
\end{exercise}

\begin{exercise}
    \(\star\star\) Un matraz A de \qty{400}{\cubic\centi\metre} contiene oxígeno a \qty{0.8}{\atm} y otro B de \qty{600}{\cubic\centi\metre}
    contiene nitrógeno a \qty{2}{\atm}, ambos a \qty{298}{\kelvin}. Se conectan los dos matraces. Si la temperatura no cambia, calcula:
    a) la presión parcial de cada gas; b) la presión total.
    \begin{solu}
        \textcolor{blue}{Cada gas pasa a ocupar \qty{1000}{\cubic\centi\metre} (Boyle):} a) \(p_{\ch{O2}} = \dfrac{\num{0.8} \cdot 400}{1000} =
        \textcolor{red}{\qty{0.32}{\atm}}\); \(p_{\ch{N2}} = \dfrac{2 \cdot 600}{1000} = \textcolor{red}{\qty{1.2}{\atm}}\). b) \(\textcolor{red}{\qty{1.52}{\atm}}\).
    \end{solu}
\end{exercise}

\begin{exercise}
    \(\star\star\) ¿Qué volumen ocupa una mezcla de \qty{10.1}{\gram} de neón y \qty{20}{\gram} de argón a \qty{3}{\atm} y \qty{27}{\degreeCelsius}?
    (Ne = 20,2; Ar = 39,9.)
    \begin{solu}
        \(n_T = \num{0.5} + \num{0.5} = \qty{1}{\mole}\); \(V = \dfrac{1 \cdot \num{0.082} \cdot 300}{3} = \textcolor{red}{\qty{8.2}{\litre}}\).
    \end{solu}
\end{exercise}

\begin{exercise}
    \(\star\star\) Un recipiente con \qty{4}{\litre} de \ch{NH3} a \qty{800}{\mmHg} se conecta con otro igual que contiene argón a
    \qty{640}{\mmHg}, a la misma temperatura. Calcula la presión total de la mezcla.
    \begin{solu}
        \textcolor{blue}{Cada gas ocupa el doble de volumen y su presión se reduce a la mitad:} \(p_T = 400 + 320 = \qty{720}{\mmHg} =
        \textcolor{red}{\qty{0.947}{\atm}}\).
    \end{solu}
\end{exercise}

%========================================================================================
\section{Ejercicios de ampliación}\label{sec:ud1-ampliacion}
%========================================================================================

\subsection{Leyes ponderales y volumétricas}

\begin{exercise}
    \(\star\) Los volcanes emiten sulfuro de hidrógeno, \ch{H2S}, que reacciona con el oxígeno del aire y forma agua y dióxido de
    azufre. Si 34 toneladas de \ch{H2S} reaccionan con 48 toneladas de oxígeno y se forman 18 toneladas de agua, ¿cuántas
    toneladas de \ch{SO2} se forman?
    \begin{solu}
        \textcolor{blue}{Lavoisier:} \(34 + 48 - 18 = \textcolor{red}{\qty{64}{\tonne}}\) de \ch{SO2}.
    \end{solu}
\end{exercise}

\begin{exercise}
    \(\star\star\) El bromuro de potasio tiene un 67,1\,\% de bromo y un 32,9\,\% de potasio. Si se mezclan \qty{20}{\gram} de bromo con
    \qty{12}{\gram} de potasio, ¿qué sustancia queda en exceso y en qué cantidad?
    \begin{solu}
        \textcolor{blue}{Para 20 g de Br hacen falta \(20 \cdot \num{32.9}/\num{67.1} = \qty{9.81}{\gram}\) de K.} \textcolor{red}{Sobran
        \(12 - \num{9.81} = \qty{2.19}{\gram}\) de potasio.}
    \end{solu}
\end{exercise}

\begin{exercise}
    \(\star\star\) Se queman en exceso de aire tres muestras de carbono puro de \qty{2.40}{\gram}, \qty{4.50}{\gram} y \qty{6.00}{\gram}, y se
    obtienen \qty{8.80}{\gram}, \qty{16.50}{\gram} y \qty{22.00}{\gram} de \ch{CO2}. a) ¿Tiene el \ch{CO2} una composición fija? b) Calcula su
    composición centesimal.
    \begin{solu}
        a) \(\num{2.40}/\num{8.80} = \num{4.50}/\num{16.50} = \num{6.00}/\num{22.00} = \num{0.273}\): \textcolor{red}{sí} (ley de Proust).
        b) \textcolor{red}{C \qty{27.3}{\percent}; O \qty{72.7}{\percent}}.
    \end{solu}
\end{exercise}

\begin{exercise}
    \(\star\star\star\) \qty{20}{\gram} de mercurio reaccionan con \qty{1.595}{\gram} de oxígeno y forman un óxido rojo. En otras condiciones,
    \qty{20}{\gram} de mercurio forman con el oxígeno un óxido negro. ¿Qué masa de oxígeno se requiere para el segundo óxido?
    a) \qty{0.798}{\gram}; b) \qty{20}{\gram}; c) \qty{3.19}{\gram}.
    \begin{solu}
        \textcolor{blue}{Por la ley de las proporciones múltiples, la masa de oxígeno ha de estar en relación sencilla con 1,595 g.
        Las opciones a y c (la mitad y el doble) cumplen la ley; la b no.} El óxido rojo es \ch{HgO} y el negro, \ch{Hg2O}:
        \textcolor{red}{a) \qty{0.798}{\gram}}.
    \end{solu}
\end{exercise}

\begin{exercise}
    \(\star\star\star\) Observa la tabla de cuatro compuestos de yodo y flúor.
    \begin{center}
        \begin{tabular}{@{}ccc@{}}
            \toprule
            \textbf{Compuesto} & \textbf{Masa de I (g)} & \textbf{Masa de F (g)} \\
            \midrule
            A & 2 & 0,2994 \\[3pt]
            B & 1 & 0,4491 \\[3pt]
            C & 0,5 & 0,3743 \\[3pt]
            D & 1,5 & 1,572 \\
            \bottomrule
        \end{tabular}
    \end{center}
    a) Demuestra que se cumple la ley de las proporciones múltiples. b) Si A es \ch{IF}, determina las fórmulas de los demás.
    \begin{solu}
        a) \textcolor{blue}{Masa de F por gramo de I:} A \num{0.1497}; B \num{0.4491}; C \num{0.7486}; D \num{1.048}. Dividiendo entre
        \num{0.1497}: \textcolor{red}{1 : 3 : 5 : 7}. b) \textcolor{red}{B: \ch{IF3}; C: \ch{IF5}; D: \ch{IF7}}.
    \end{solu}
\end{exercise}

\begin{exercise}
    \(\star\star\star\) En una muestra hay cuatro óxidos de cloro, de fórmulas \ch{Cl2O}, \ch{Cl2O3}, \ch{Cl2O5} y \ch{Cl2O7}. Se analizan y
    se obtiene: el primero tiene un 18,39\,\% de oxígeno; el segundo, un 40,34\,\%; en el tercero, \qty{20}{\gram} de oxígeno reaccionan con
    \qty{17.75}{\gram} de cloro; en el cuarto, \qty{20}{\gram} de oxígeno con \qty{10}{\gram} de cloro. Uno de los análisis es incorrecto:
    ¿cuál?
    \begin{solu}
        \textcolor{blue}{Masa de O por gramo de Cl:} 1.º \(\num{18.39}/\num{81.61} = \num{0.225}\); 2.º \(\num{40.34}/\num{59.66} = \num{0.676}\);
        3.º \(20/\num{17.75} = \num{1.127}\); 4.º \(20/10 = 2\). Dividiendo entre \num{0.225}: 1 : 3 : 5 : 8,9.
        \textcolor{red}{El cuarto es incorrecto}: \textcolor{blue}{para el \ch{Cl2O7} saldría 7 (\qty{1.578}{\gram} de O por gramo de Cl).}
    \end{solu}
\end{exercise}

\begin{exercise}
    \(\star\star\) Demuestra que se cumple la ley de Dalton en dos óxidos de nitrógeno que tienen un 46,67\,\% y un 36,84\,\% de
    nitrógeno, respectivamente.
    \begin{solu}
        O por gramo de N: \(\num{53.33}/\num{46.67} = \num{1.143}\) y \(\num{63.16}/\num{36.84} = \num{1.714}\);
        \(\num{1.714}/\num{1.143} = \num{1.5} = \textcolor{red}{3/2}\). \textcolor{red}{Se cumple} (son \ch{NO} y \ch{N2O3}).
    \end{solu}
\end{exercise}

\begin{exercise}
    \(\star\star\) El azufre forma distintos compuestos con el oxígeno. En el primero, \qty{2}{\gram} de S se combinan con \qty{1.996}{\gram}
    de O; en el segundo, \qty{3}{\gram} de S con \qty{4.491}{\gram} de O. ¿Son estos resultados coherentes con la ley de Dalton?
    \begin{solu}
        O por gramo de S: \num{0.998} y \num{1.497}; \(\num{1.497}/\num{0.998} = \num{1.5}\): \textcolor{red}{sí, relación 2:3} (\ch{SO2} y \ch{SO3}).
    \end{solu}
\end{exercise}

\begin{exercise}
    \(\star\star\star\) El hidrógeno y el flúor reaccionan para dar fluoruro de hidrógeno. Demuestra con los datos de la tabla
    (volúmenes en litros, en las mismas condiciones) que se cumple la ley de Gay-Lussac.
    \begin{center}
        \begin{tabular}{@{}ccccc@{}}
            \toprule
            \multicolumn{2}{c}{\textbf{Inicial}} & \multicolumn{3}{c}{\textbf{Final}} \\
            \cmidrule(r){1-2} \cmidrule(l){3-5}
            \ch{H2} & \ch{F2} & \ch{HF} & \ch{H2} & \ch{F2} \\
            \midrule
            40 & 55 & 80 & 0 & 15 \\[3pt]
            65 & 25 & 50 & 40 & 0 \\
            \bottomrule
        \end{tabular}
    \end{center}
    \begin{solu}
        Exp. 1: reaccionan 40 L de \ch{H2} con 40 L de \ch{F2} y dan 80 L de HF. Exp. 2: reaccionan 25 L de \ch{H2} con 25 L de \ch{F2}
        y dan 50 L. \textcolor{red}{En ambas, la relación es 1 : 1 : 2}: números enteros sencillos.
    \end{solu}
\end{exercise}

\begin{exercise}
    \(\star\star\) El fósforo forma dos cloruros: \ch{PCl3} y \ch{PCl5}. a) Calcula la relación entre las masas de Cl y P en cada
    uno. b) Si en el primero hay \qty{3.435}{\gram} de cloro por cada gramo de fósforo, ¿cuánto cloro se combina con un gramo de
    fósforo en el segundo?
    \begin{solu}
        a) \ch{PCl3}: \(\num{106.5}/31 = \textcolor{red}{\num{3.435}}\); \ch{PCl5}: \(\num{177.5}/31 = \textcolor{red}{\num{5.726}}\).
        b) \(\num{3.435} \cdot 5/3 = \textcolor{red}{\qty{5.73}{\gram}}\).
    \end{solu}
\end{exercise}

\subsection{Cantidad de sustancia y fórmulas}

\begin{exercise}
    \(\star\) El estireno, materia prima del poliestireno, es un hidrocarburo de fórmula molecular \ch{C8H8}. Calcula: a) su fórmula
    empírica; b) las moléculas que hay en \qty{52}{\gram}.
    \begin{solu}
        a) \textcolor{red}{\ch{CH}}. b) \(M = 104\); \(n = \qty{0.5}{\mole}\); \(N = \textcolor{red}{\num{3.01e23}}\) moléculas.
    \end{solu}
\end{exercise}

\begin{exercise}
    \(\star\star\star\) El oro se vende a unos 2000 dólares la onza. Si el dólar se cambia a 0,92 euros y una onza son \qty{31.1}{\gram},
    ¿cuánto cuesta un mol de oro? (Au = 196,97.)
    \begin{solu}
        \(\dfrac{\num{196.97}}{\num{31.1}} \cdot 2000 \cdot \num{0.92} = \textcolor{red}{\num{11650}\text{ €}}\).
    \end{solu}
\end{exercise}

\begin{exercise}
    \(\star\star\star\) El fosfato de cinc se usa como cemento dental. Una muestra de \qty{40.0}{\milli\gram} se descompone en sus
    elementos y se obtienen \qty{13.26}{\milli\gram} de oxígeno, \qty{6.42}{\milli\gram} de fósforo y \qty{20.32}{\milli\gram} de cinc.
    Determina la fórmula empírica. (Zn = 65,4; P = 31.)
    \begin{solu}
        Zn \(\num{20.32}/\num{65.4} = \num{0.311}\); P \(\num{6.42}/31 = \num{0.207}\); O \(\num{13.26}/16 = \num{0.829}\) \(\rightarrow\) 1,5 : 1 : 4
        \(\rightarrow\) (×2) \textcolor{red}{\ch{Zn3P2O8}}, es decir, \ch{Zn3(PO4)2}.
    \end{solu}
\end{exercise}

\begin{exercise}
    \(\star\star\) La masa atómica del sodio es \qty{23}{\uma}. Calcula: a) la masa en gramos de un átomo de Na; b) cuál de estas
    cantidades tiene más átomos: \qty{46}{\gram} de Na, \qty{1.5}{\mole} de Na o \num{1e24} átomos de Na.
    \begin{solu}
        a) \(23/\num{6.022e23} = \textcolor{red}{\qty{3.82e-23}{\gram}}\). b) \qty{46}{\gram} son \qty{2}{\mole}; \num{1e24} átomos son \qty{1.66}{\mole}:
        \textcolor{red}{\qty{46}{\gram} de Na}.
    \end{solu}
\end{exercise}

\begin{exercise}
    \(\star\star\) La masa atómica del cobre es \qty{63.5}{\uma}. a) ¿Cuántos átomos hay en \qty{2}{\gram} de cobre? b) ¿Cuál de estas cantidades
    tiene mayor masa: \qty{63.5}{\gram} de Cu, \qty{0.8}{\mole} de Cu o \num{7e23} átomos de Cu?
    \begin{solu}
        a) \(\dfrac{2}{\num{63.5}} \cdot \num{6.022e23} = \textcolor{red}{\num{1.90e22}}\). b) \qty{0.8}{\mole} son \qty{50.8}{\gram}; \num{7e23} átomos son
        \qty{73.8}{\gram}: \textcolor{red}{los \num{7e23} átomos}.
    \end{solu}
\end{exercise}

\begin{exercise}
    \(\star\star\star\) La morfina, \ch{C17H19NO3}, se usa para tratar dolores intensos. La dosis diaria de un paciente es de
    \qty{15}{\milli\gram}. a) ¿Cuántos moles son? b) ¿Cuántas moléculas? c) ¿Cuántos átomos de C y de O? d) Si se reparte en tomas cada
    8 horas, ¿qué cantidad hay en cada toma? ¿Cuántos moles?
    \begin{solu}
        \(M = \qty{285}{\gram\per\mole}\). a) \(\num{0.015}/285 = \textcolor{red}{\qty{5.26e-5}{\mole}}\). b) \(\textcolor{red}{\num{3.17e19}}\).
        c) C: \(17 \cdot \num{3.17e19} = \textcolor{red}{\num{5.39e20}}\); O: \(3 \cdot \num{3.17e19} = \textcolor{red}{\num{9.51e19}}\).
        d) 3 tomas al día: \textcolor{red}{\qty{5}{\milli\gram}} por toma, \textcolor{red}{\qty{1.75e-5}{\mole}}.
    \end{solu}
\end{exercise}

\begin{exercise}
    \(\star\star\) Halla: a) el número de átomos de neón que hay en \qty{80}{\milli\gram} de neón; b) la masa de oxígeno, \ch{O2}, que
    tiene el mismo número de moléculas que \qty{142}{\gram} de cloro, \ch{Cl2}. (Ne = 20,2.)
    \begin{solu}
        a) \(\dfrac{\num{0.08}}{\num{20.2}} \cdot \num{6.022e23} = \textcolor{red}{\num{2.38e21}}\). b) \(142/71 = \qty{2}{\mole}\) \(\rightarrow\)
        \(2 \cdot 32 = \textcolor{red}{\qty{64}{\gram}}\).
    \end{solu}
\end{exercise}

\begin{exercise}
    \(\star\) Calcula el número de átomos de Fe que hay en una pieza de \qty{279}{\gram} de hierro.
    \begin{solu}
        \(279/\num{55.8} = \qty{5}{\mole}\) \(\rightarrow\) \(\textcolor{red}{\num{3.01e24}}\) átomos.
    \end{solu}
\end{exercise}

\begin{exercise}
    \(\star\star\star\) Una dosis de aspirina, \ch{C9H8O4}, contiene \qty{6e-4}{\mole}. ¿Cuántas dosis pueden prepararse con
    \qty{3}{\kilogram} de aspirina?
    \begin{solu}
        \(M = 180\); \(n = 3000/180 = \qty{16.7}{\mole}\); \(\num{16.7}/\num{6e-4} = \textcolor{red}{\num{27800}}\) dosis.
    \end{solu}
\end{exercise}

\begin{exercise}
    \(\star\star\star\) En una ciudad de \num{500000} habitantes se produce una epidemia de una infección cutánea. Cada paciente debe tomar 2
    pastillas de ibuprofeno al día durante 10 días, y se calcula que hacen falta 2 toneladas de ibuprofeno, \ch{C13H18O2}, para tratar a
    toda la población. a) ¿Qué masa tiene cada pastilla? b) ¿Qué masa toma cada paciente al día? c) ¿Cuántos átomos hay en un miligramo de
    ibuprofeno? d) ¿Cuántas moléculas toma un paciente al día?
    \begin{solu}
        a) \(\dfrac{\num{2e6}\text{ g}}{\num{500000} \cdot 20} = \textcolor{red}{\qty{200}{\milli\gram}}\). b) \textcolor{red}{\qty{400}{\milli\gram}}.
        c) \(M = 206\); 33 átomos por molécula: \(\dfrac{\num{1e-3}}{206} \cdot \num{6.022e23} \cdot 33 = \textcolor{red}{\num{9.65e19}}\).
        d) \(\dfrac{\num{0.4}}{206} \cdot \num{6.022e23} = \textcolor{red}{\num{1.17e21}}\).
    \end{solu}
\end{exercise}

\begin{exercise}
    \(\star\star\star\) Algunos peces contienen plomo. Supón que el máximo permitido en sangre es de \qty{200}{\micro\gram} de Pb por litro. Determina:
    a) el máximo en \unit{\gram\per\litre}; b) en \unit{\mole\per\litre}; c) el número de átomos de plomo por litro de sangre.
    \begin{solu}
        a) \(\textcolor{red}{\qty{2e-4}{\gram\per\litre}}\). b) \(\num{2e-4}/\num{207.2} = \textcolor{red}{\qty{9.65e-7}{\mole\per\litre}}\).
        c) \(\textcolor{red}{\num{5.81e17}}\) átomos.
    \end{solu}
\end{exercise}

\subsection{Leyes de los gases}

\begin{exercise}
    \(\star\star\) Una rueda puede explotar si la presión en su interior supera la máxima permitida. Explica cómo podemos evitar que
    explote el neumático si se duplica su presión.
    \begin{solu}
        \textcolor{red}{Liberando aire} (menos moles, menos presión) o \textcolor{red}{enfriándolo}: \textcolor{blue}{a volumen casi constante,
        \(p\) es proporcional a \(n\) y a \(T\).}
    \end{solu}
\end{exercise}

\begin{exercise}
    \(\star\star\) Una muestra de oxígeno ocupa \qty{60}{\milli\litre} a la presión atmosférica. Si la presión aumenta en
    \qty{100}{\mmHg}, el volumen se reduce a \qty{48}{\milli\litre}. ¿Cuál era la presión atmosférica, en \unit{\mmHg} y en \unit{\atm}?
    \begin{solu}
        \(p \cdot 60 = (p + 100) \cdot 48 \Rightarrow 12\,p = 4800 \Rightarrow p = \textcolor{red}{\qty{400}{\mmHg}} = \textcolor{red}{\qty{0.526}{\atm}}\).
        \textcolor{blue}{(Una presión tan baja corresponde a una gran altitud.)}
    \end{solu}
\end{exercise}

\begin{exercise}
    \(\star\star\) Un globo lleno de helio ocupa \qty{1500}{\cubic\metre} al nivel del suelo, a \qty{1}{\atm} y \qty{20}{\degreeCelsius}. Al
    ascender hasta donde la presión es \qty{0.5}{\atm}, su volumen pasa a \qty{2600}{\cubic\metre}. ¿Cuál es la temperatura a esa altitud?
    \begin{solu}
        \(T_2 = T_1 \dfrac{p_2 V_2}{p_1 V_1} = 293 \cdot \dfrac{\num{0.5} \cdot 2600}{1 \cdot 1500} = \textcolor{red}{\qty{254}{\kelvin}}\) (\qty{-19}{\degreeCelsius}).
    \end{solu}
\end{exercise}

\begin{exercise}
    \(\star\star\) Un recipiente de \qty{4}{\litre} soporta como máximo \qty{3e5}{\pascal}. Si contiene \qty{16}{\gram} de \ch{SO2}, ¿hasta qué
    temperatura puede calentarse?
    \begin{solu}
        \(n = \qty{0.25}{\mole}\); \(p = \num{3e5}/\num{1.013e5} = \qty{2.96}{\atm}\); \(T = \dfrac{\num{2.96} \cdot 4}{\num{0.25} \cdot \num{0.082}} =
        \textcolor{red}{\qty{578}{\kelvin}}\).
    \end{solu}
\end{exercise}

\begin{exercise}
    \(\star\star\) Queremos aumentar el volumen de un gas desde \qty{500}{\cubic\centi\metre} hasta \qty{2}{\litre} a temperatura constante. Si la
    presión inicial es de \qty{2e5}{\pascal}, ¿qué presión final debe alcanzar?
    \begin{solu}
        \(p_2 = \dfrac{\num{2e5} \cdot \num{0.5}}{2} = \textcolor{red}{\qty{5e4}{\pascal}} = \textcolor{red}{\qty{0.49}{\atm}}\).
    \end{solu}
\end{exercise}

\begin{exercise}
    \(\star\star\star\) ¿Por qué hay que medir la presión de los neumáticos con el coche parado y sin haber recorrido grandes distancias?
    Explica qué ley justifica esta recomendación.
    \begin{solu}
        \textcolor{blue}{Al rodar, el rozamiento calienta los neumáticos y el aire de su interior.} \textcolor{red}{Por la ley de Gay-Lussac, al
        aumentar la temperatura a volumen casi constante, aumenta la presión}: mediríamos un valor mayor que el real en frío, que es el
        de referencia.
    \end{solu}
\end{exercise}

\begin{exercise}
    \(\star\star\star\) Luis compró una bolsa de patatas fritas en una playa de Tenerife y al día siguiente subió con ella al Teide. Al
    subir, observó que la bolsa se hinchaba. Explica lo sucedido.
    \begin{solu}
        \textcolor{blue}{La bolsa está cerrada con aire a la presión del nivel del mar.} \textcolor{red}{Al subir, la presión exterior disminuye
        y, por la ley de Boyle, el aire de la bolsa se expande hasta que su presión iguala a la de fuera}: la bolsa se hincha (aunque la
        temperatura baje algo con la altura, domina la disminución de presión).
    \end{solu}
\end{exercise}

\begin{exercise}
    \(\star\star\star\) Los buceadores saben que no deben ascender deprisa desde gran profundidad, y que deben exhalar aire de forma
    constante al subir. Explica qué le pasaría a los pulmones de un buceador que tomara aire a presión a 30 m de profundidad y
    subiera a la superficie conteniendo la respiración.
    \begin{solu}
        \textcolor{blue}{A 30 m la presión es de unas 4 atm (1 de la atmósfera y unas 3 del agua).} \textcolor{red}{Al subir, la presión baja
        a 1 atm y, por la ley de Boyle, el aire de los pulmones se expandiría hasta unas 4 veces su volumen}, lo que podría dañarlos
        gravemente. Además, los gases disueltos en la sangre a alta presión formarían burbujas si se asciende demasiado rápido.
    \end{solu}
\end{exercise}

\subsection{Ecuación de los gases ideales}

\begin{exercise}
    \(\star\star\) ¿Cuántas bombonas de oxígeno de \qty{5}{\litre} se pueden llenar a \qty{1013}{\milli\bar} y \qty{15}{\degreeCelsius} con un
    tanque de \qty{4000}{\litre} que contiene oxígeno a \qty{6e5}{\pascal} y \qty{20}{\degreeCelsius}?
    \begin{solu}
        Tanque: \(n = \dfrac{\num{6e5} \cdot 4}{\num{8.31} \cdot 293} = \qty{986}{\mole}\). Bombona: \(n = \dfrac{\num{1.013e5} \cdot \num{0.005}}{\num{8.31} \cdot 288} =
        \qty{0.212}{\mole}\). \(986/\num{0.212} = \textcolor{red}{\num{4657}}\) \textcolor{red}{bombonas}.
    \end{solu}
\end{exercise}

\begin{exercise}
    \(\star\star\star\) Los puntos A, B, C y D de la gráfica representan cuatro estados de una misma muestra de gas, situados sobre dos
    isotermas (curvas a temperatura constante) \(T_1\) y \(T_2\).
    \begin{center}
        \begin{tikzpicture}
            \begin{axis}[grafica, width=0.55\linewidth, xlabel={\(V\) (\unit{\litre})}, ylabel={\(p\) (\unit{\atm})},
                xmin=0, xmax=4.5, ymin=0, ymax=6.5]
                \addplot[thick, green!50!black, domain=0.5:4.4] {3/x};
                \addplot[thick, blue, domain=1:4.4] {6/x};
                \node[font=\scriptsize, blue, anchor=south] at (axis cs:4.2,1.45) {\(T_2\)};
                \node[font=\scriptsize, green!50!black, anchor=north] at (axis cs:4.2,0.7) {\(T_1\)};
                \addplot[only marks, mark=*, red] coordinates {(1,3) (2,3) (3,2) (2,1.5)};
                \node[font=\scriptsize, anchor=south east] at (axis cs:1,3) {A};
                \node[font=\scriptsize, anchor=south west] at (axis cs:2,3) {B};
                \node[font=\scriptsize, anchor=south west] at (axis cs:3,2) {C};
                \node[font=\scriptsize, anchor=north east] at (axis cs:2,1.5) {D};
            \end{axis}
        \end{tikzpicture}
    \end{center}
    a) Indica dos estados con la misma temperatura. b) ¿En qué estado la presión es la misma que en A? c) ¿En qué estados el volumen es
    el mismo? d) Explica qué pasa con \(p\), \(V\) y \(T\) en los procesos AB, BC y CD. e) ¿Qué relación hay entre \(T_2\) y \(T_1\)?
    \begin{solu}
        a) \textcolor{red}{A y D (\(T_1\)); B y C (\(T_2\))}. b) \textcolor{red}{B}. c) \textcolor{red}{B y D} (\qty{2}{\litre}).
        d) \textcolor{red}{AB}: \(p\) constante, \(V\) y \(T\) aumentan. \textcolor{red}{BC}: \(T\) constante, \(V\) aumenta y \(p\) disminuye.
        \textcolor{red}{CD}: \(p\), \(V\) y \(T\) disminuyen. e) \textcolor{blue}{\(pV\) es proporcional a \(T\): \(pV = 6\) en \(T_2\) y \(3\) en \(T_1\).}
        \(\textcolor{red}{T_2/T_1 = 2}\).
    \end{solu}
\end{exercise}

\begin{exercise}
    \(\star\star\) Halla la presión que ejercen \qty{100}{\gram} de \ch{CO2} en un recipiente cilíndrico de \qty{30}{\centi\metre} de diámetro y
    \qty{1}{\metre} de altura a \qty{10}{\degreeCelsius}.
    \begin{solu}
        \(V = \pi \cdot \num{0.15}^2 \cdot 1 = \qty{0.0707}{\cubic\metre} = \qty{70.7}{\litre}\); \(n = \qty{2.27}{\mole}\);
        \(p = \dfrac{\num{2.27} \cdot \num{0.082} \cdot 283}{\num{70.7}} = \textcolor{red}{\qty{0.746}{\atm}}\).
    \end{solu}
\end{exercise}

\begin{exercise}
    \(\star\star\star\) Una persona expulsa unos \qty{35}{\gram} de dióxido de carbono por hora. Calcula el volumen de \ch{CO2}, a
    \qty{37}{\degreeCelsius} y \qty{1}{\atm}, que expulsa al día.
    \begin{solu}
        \(n = \dfrac{35 \cdot 24}{44} = \qty{19.1}{\mole}\); \(V = \num{19.1} \cdot \num{0.082} \cdot 310 = \textcolor{red}{\qty{485}{\litre}}\).
    \end{solu}
\end{exercise}

\begin{exercise}
    \(\star\star\star\) Sabiendo que \qty{1}{\litre} de aire en condiciones normales tiene una masa de \qty{1.293}{\gram}, calcula la masa de
    aire contenida en un recipiente de \qty{2500}{\cubic\centi\metre} a \qty{10}{\atm} y \qty{20}{\degreeCelsius}.
    \begin{solu}
        \(M_{\text{aire}} = \num{1.293} \cdot \num{22.4} = \qty{28.96}{\gram\per\mole}\); \(n = \dfrac{10 \cdot \num{2.5}}{\num{0.082} \cdot 293} =
        \qty{1.04}{\mole}\); \(m = \textcolor{red}{\qty{30.1}{\gram}}\).
    \end{solu}
\end{exercise}

\begin{exercise}
    \(\star\star\) Un tanque de \qty{4000}{\cubic\centi\metre} contiene oxígeno a \qty{25}{\degreeCelsius}. Su manómetro marca \qty{500}{\kilo\pascal}
    (los manómetros miden la diferencia con la presión atmosférica, que es de \qty{101}{\kilo\pascal}). ¿Qué masa de oxígeno hay?
    \begin{solu}
        \(p = 500 + 101 = \qty{601}{\kilo\pascal}\). \(n = \dfrac{\num{6.01e5} \cdot \num{0.004}}{\num{8.31} \cdot 298} = \qty{0.971}{\mole}\);
        \(m = \textcolor{red}{\qty{31.1}{\gram}}\).
    \end{solu}
\end{exercise}

\begin{exercise}
    \(\star\star\star\) Un hidrocarburo gaseoso contiene un 81,8\,\% de carbono y un 18,2\,\% de hidrógeno. \qty{6.6}{\gram} de él en una
    ampolla de \qty{2}{\litre} a \qty{50}{\degreeCelsius} ejercen una presión de \qty{1.99}{\atm}. Determina su fórmula molecular.
    \begin{solu}
        C \num{6.82}; H \num{18.2} \(\rightarrow\) 1 : 2,67 \(\rightarrow\) (×3) \ch{C3H8} (\(M = 44\)).
        \(M = \dfrac{\num{6.6} \cdot \num{0.082} \cdot 323}{\num{1.99} \cdot 2} = \qty{43.9}{\gram\per\mole}\); \(\alpha = 1\): \textcolor{red}{\ch{C3H8}} (propano).
    \end{solu}
\end{exercise}

\begin{exercise}
    \(\star\star\star\) Un compuesto orgánico contiene un 54,5\,\% de C, un 36,4\,\% de O y un 9,1\,\% de H. \qty{3.5}{\gram} de su vapor ocupan
    \qty{1.26}{\litre} a \qty{750}{\mmHg} y \qty{380}{\kelvin}. Determina su fórmula molecular.
    \begin{solu}
        C \num{4.54}; H \num{9.1}; O \num{2.275} \(\rightarrow\) 2 : 4 : 1: \ch{C2H4O} (\(M = 44\)).
        \(M = \dfrac{\num{3.5} \cdot \num{0.082} \cdot 380}{\num{0.987} \cdot \num{1.26}} = \qty{87.7}{\gram\per\mole}\); \(\alpha = 2\): \textcolor{red}{\ch{C4H8O2}}.
    \end{solu}
\end{exercise}

\subsection{Densidad de los gases}

\begin{exercise}
    \(\star\star\) El cloroeteno (cloruro de vinilo) se usa para fabricar PVC. Su densidad es \qty{2.8}{\gram\per\litre} a \qty{25}{\degreeCelsius} y
    \qty{1.1}{\atm}. Calcula su masa molar.
    \begin{solu}
        \(M = \dfrac{d R T}{p} = \dfrac{\num{2.8} \cdot \num{0.082} \cdot 298}{\num{1.1}} = \textcolor{red}{\qty{62.2}{\gram\per\mole}}\) (\ch{C2H3Cl}).
    \end{solu}
\end{exercise}

\begin{exercise}
    \(\star\star\) ¿Cuál de estos gases tiene mayor densidad en condiciones normales: a) \ch{NH3}; b) \ch{N2}; c) \ch{SO2}; d) Ar? ¿Y a
    \qty{25}{\degreeCelsius} y \qty{3}{\atm}?
    \begin{solu}
        \textcolor{blue}{A igual \(p\) y \(T\), \(d\) es proporcional a \(M\):} \textcolor{red}{c) \ch{SO2}} (\(M = 64\)) en ambos casos.
    \end{solu}
\end{exercise}

\begin{exercise}
    \(\star\star\star\) A un niño se le escapa un globo de helio y sube volando. Para que un globo suba, la densidad del gas que contiene
    debe ser menor que la del aire. a) Calcula la densidad del aire a \qty{1}{\atm} y \qty{15}{\degreeCelsius} (masa molar media
    \qty{29}{\gram\per\mole}). b) Demuestra que un globo lleno de \ch{CO2} en esas condiciones no podría volar. c) ¿A qué temperatura
    habría que calentar el \ch{CO2} para que su densidad igualara la del aire a \qty{15}{\degreeCelsius}?
    \begin{solu}
        a) \(d = \dfrac{1 \cdot 29}{\num{0.082} \cdot 288} = \textcolor{red}{\qty{1.23}{\gram\per\litre}}\). b) \(d_{\ch{CO2}} = \dfrac{44}{\num{0.082} \cdot 288} =
        \qty{1.86}{\gram\per\litre}\), \textcolor{red}{mayor que la del aire}. c) \(T = \dfrac{p M}{d R} = \dfrac{44}{\num{1.23} \cdot \num{0.082}} =
        \textcolor{red}{\qty{437}{\kelvin}}\).
    \end{solu}
\end{exercise}

\begin{exercise}
    \(\star\star\) El propeno, \ch{C3H6}, se usa para fabricar polipropileno. Un matraz de \qty{150}{\milli\litre} tiene una masa de
    \qty{52.3410}{\gram} vacío. ¿Qué masa tendrá lleno de propeno a \qty{750}{\mmHg} y \qty{20}{\degreeCelsius}?
    \begin{solu}
        \(n = \dfrac{\num{0.987} \cdot \num{0.15}}{\num{0.082} \cdot 293} = \qty{6.16e-3}{\mole}\); \(m = \num{6.16e-3} \cdot 42 = \qty{0.2588}{\gram}\).
        \textcolor{red}{Masa total: \qty{52.5998}{\gram}}.
    \end{solu}
\end{exercise}

\begin{exercise}
    \(\star\star\star\) La temperatura del centro del Sol se puede estimar suponiendo que su núcleo es un gas ideal de masa molar media
    \qty{0.6}{\gram\per\mole}, a una presión de \qty{2.4e11}{\atm} y con una densidad de \qty{1.5e5}{\kilogram\per\cubic\metre}. Calcula esa
    temperatura.
    \begin{solu}
        \(d = \qty{1.5e5}{\gram\per\litre}\). \(T = \dfrac{p M}{d R} = \dfrac{\num{2.4e11} \cdot \num{0.6}}{\num{1.5e5} \cdot \num{0.082}} =
        \textcolor{red}{\qty{1.2e7}{\kelvin}}\).
    \end{solu}
\end{exercise}

\begin{exercise}
    \(\star\star\star\) Del fondo de un lago de \qty{15}{\metre} de profundidad sale una burbuja de aire de \qty{8}{\milli\litre}. El agua del fondo está
    a \qty{277}{\kelvin}, y la de la superficie, a \qty{15}{\degreeCelsius}. ¿Qué volumen tendrá la burbuja al llegar a la superficie?
    La presión atmosférica es de \qty{101325}{\pascal}; la presión del agua es \(p = d\,g\,h\), con \(d = \qty{1000}{\kilogram\per\cubic\metre}\) y
    \(g = \qty{9.8}{\metre\per\second\squared}\).
    \begin{solu}
        \(p_1 = 101325 + 1000 \cdot \num{9.8} \cdot 15 = \qty{248325}{\pascal}\).
        \(V_2 = V_1 \dfrac{p_1 T_2}{p_2 T_1} = 8 \cdot \dfrac{248325 \cdot 288}{101325 \cdot 277} = \textcolor{red}{\qty{20.4}{\milli\litre}}\).
    \end{solu}
\end{exercise}

\subsection{Mezclas de gases y presiones parciales}

\begin{exercise}
    \(\star\star\) Una bombona de \qty{25}{\litre} a \qty{20}{\degreeCelsius} lleva un manómetro y contiene \qty{56}{\gram} de nitrógeno,
    \qty{12}{\gram} de helio y \qty{22}{\gram} de dióxido de carbono. Calcula: a) qué gas ejerce mayor presión; b) la presión total; c) qué
    marcará el manómetro si mide la diferencia con la presión atmosférica (\qty{1}{\atm}).
    \begin{solu}
        \(n\): \ch{N2} 2; He 3; \ch{CO2} \num{0.5}. a) \textcolor{red}{El helio} (más moles): \(p_{\ch{He}} = \qty{2.88}{\atm}\);
        \(p_{\ch{N2}} = \qty{1.92}{\atm}\); \(p_{\ch{CO2}} = \qty{0.48}{\atm}\). b) \(\textcolor{red}{\qty{5.29}{\atm}}\). c) \textcolor{red}{\qty{4.29}{\atm}}.
    \end{solu}
\end{exercise}

\begin{exercise}
    \(\star\star\) Una mezcla gaseosa a \qty{80}{\degreeCelsius} está formada por \qty{255}{\gram} de amoníaco y una masa desconocida de dióxido de
    carbono. Las presiones parciales son \qty{2e5}{\pascal} para el amoníaco y \qty{4e4}{\pascal} para el \ch{CO2}. Halla: a) la presión total;
    b) la masa de \ch{CO2}; c) el volumen de la mezcla.
    \begin{solu}
        a) \(\textcolor{red}{\qty{2.4e5}{\pascal}}\). b) \(n_{\ch{NH3}} = \qty{15}{\mole}\); \(n_{\ch{CO2}} = 15 \cdot \dfrac{\num{4e4}}{\num{2e5}} = \qty{3}{\mole}\):
        \(\textcolor{red}{\qty{132}{\gram}}\). c) \(V = \dfrac{15 \cdot \num{8.31} \cdot 353}{\num{2e5}} = \qty{0.22}{\cubic\metre} = \textcolor{red}{\qty{220}{\litre}}\).
    \end{solu}
\end{exercise}

\begin{exercise}
    \(\star\star\) Un recipiente A contiene \qty{34}{\gram} de amoníaco y otro B, \qty{32}{\gram} de metano. Ambos tienen \qty{30}{\litre} y están a
    \qty{100}{\degreeCelsius}. Los dos gases se pasan a un tercer recipiente C de \qty{80}{\litre} a \qty{25}{\degreeCelsius}. Calcula la presión
    en A, en B y en C.
    \begin{solu}
        \(n_A = n_B = \qty{2}{\mole}\). \(p_A = p_B = \dfrac{2 \cdot \num{0.082} \cdot 373}{30} = \textcolor{red}{\qty{2.04}{\atm}}\);
        \(p_C = \dfrac{4 \cdot \num{0.082} \cdot 298}{80} = \textcolor{red}{\qty{1.22}{\atm}}\).
    \end{solu}
\end{exercise}

\begin{exercise}
    \(\star\star\) Una botella de \qty{3}{\litre} contiene \qty{3.2}{\gram} de oxígeno a \qty{273}{\kelvin}. ¿Cuál de estas opciones permitiría
    que la presión fuera de \qty{2}{\atm}, sin cambiar la temperatura? a) Añadir \qty{3.2}{\gram} de oxígeno. b) Añadir \qty{5.4}{\gram} de
    oxígeno. c) Añadir \qty{2}{\gram} de helio. d) Liberar \qty{1.6}{\gram} de oxígeno.
    \begin{solu}
        \textcolor{blue}{Moles necesarios: \(n = \dfrac{2 \cdot 3}{\num{0.082} \cdot 273} = \qty{0.268}{\mole}\); hay \qty{0.1}{\mole}, faltan
        \qty{0.168}{\mole}.} a) añade 0,1 mol; \textcolor{red}{b) añade \(\num{5.4}/32 = \qty{0.169}{\mole}\): es la correcta}; c) añade 0,5 mol;
        d) quita moles.
    \end{solu}
\end{exercise}

\begin{exercise}
    \(\star\star\star\) El aire espirado es distinto del aire normal. Un análisis del aire espirado a \qty{37}{\degreeCelsius} y \qty{1}{\atm} da, en
    volumen: 75\,\% de nitrógeno, 16\,\% de oxígeno, 4\,\% de dióxido de carbono, 4\,\% de vapor de agua y 1\,\% de argón. Determina:
    a) las presiones parciales del nitrógeno y del oxígeno en el aire espirado y en el aire normal (78\,\% de \ch{N2}, 21\,\% de \ch{O2},
    1\,\% de Ar); b) si la densidad del aire espirado es mayor o menor que la del aire normal, en las mismas condiciones;
    c) las presiones parciales del apartado a si el aire espirado estuviera a \qty{2}{\atm}.
    \begin{solu}
        a) Espirado: \(p_{\ch{N2}} = \textcolor{red}{\qty{0.75}{\atm}}\), \(p_{\ch{O2}} = \textcolor{red}{\qty{0.16}{\atm}}\); normal: \textcolor{red}{\qty{0.78}{\atm}} y
        \textcolor{red}{\qty{0.21}{\atm}}.
        b) \textcolor{blue}{Masa molar media:} espirado \(\num{0.75} \cdot 28 + \num{0.16} \cdot 32 + \num{0.04} \cdot 44 + \num{0.04} \cdot 18 + \num{0.01} \cdot \num{39.9} =
        \qty{29.00}{\gram\per\mole}\); normal \qty{28.96}{\gram\per\mole}: \textcolor{red}{el espirado es ligeramente más denso}.
        c) \textcolor{red}{El doble: \qty{1.5}{\atm} y \qty{0.32}{\atm}}.
    \end{solu}
\end{exercise}
